\documentclass[11pt]{article}
\usepackage[T1]{fontenc}
\usepackage{lmodern,amsmath,amssymb,amsthm,booktabs,tabularx,hyperref}
\usepackage[margin=1in]{geometry}
\newtheorem{theorem}{Theorem}
\newcommand{\Sol}{\operatorname{Sol}}
\newcommand{\Models}{\operatorname{Models}}
\newcommand{\code}[1]{\texttt{\small #1}}
\title{A Transform Candidate for Higher-Order MeTTa Prime\\
\large Atomic interaction, contextual solutions and justified absence}
\author{Oru\v{z}i}
\date{5 November 2026}
\begin{document}
\maketitle

\begin{abstract}
The candidate treats communication, query selection and rule execution as
uses of one guarded occurrence transformation. Its information component is
the joint solution set of constraints, not one chosen unifier. We describe
the existing atomic rule, prove a sufficient guarded commutation law with
nontrivial shared equations, compare actual contextual solver outputs across
different helper declarations, characterize coverage sufficient for
negative queries, and compare finite clause programs with actual solver
histories. These results support a common interaction mechanism.
They do not yet establish a compositional foundation for every required
Prime program or an unrestricted equivalence with other MeTTa semantics.
\end{abstract}

\section{The question and the comparison criterion}

An adequate common calculus must explain how programs are formed and how
their behavior is preserved in surrounding programs. Calling an opaque
evaluator from a transformation does not establish this correspondence.
The comparison sources remain independently specified: the register
semantics of Meta-MeTTa~\cite{metametta}, HE's minimal instructions, the
declared rho-calculus profile, the symmetric reflective MeTTa-calculus,
and the Prime proposals.

An exact restriction, a compositional encoding and an essential extension
are different possible results. Each comparison must declare its source
fragment and observations: answers, repeated occurrences, required order,
binding correlations, effects and suspension. A comparison for translated
source contexts is not automatically valid for all richer target contexts.
An inability to complete a proof is not an impossibility result.

The first independent register comparison already separates a local
prototype from the published source. The finite extensional source profile
in \code{OSLFCore.MetaMeTTa} records the eight displayed register rules,
a declared principal-unifier interpretation and a named choice about empty match families.
Comprehensions, streams and host arithmetic are outside this profile.
With the same input representation and exact output observation, the older
workspace prototype cannot simulate one Transform example even through
arbitrarily many finite steps. A distinct encoding is not excluded.

Several boundaries matter for the common calculus. A source answer
family can use $x=a$ and $x=b$ in separate instances; conjoining them under
one shared variable loses those answers. A concrete CMR comparison rejects
that identification and constructs a two-answer firing using fresh private
variables, preserving the caller's observations outside those names.
This selected example is not a uniform compiler: unrestricted core reads
also permit duplicate selection of one knowledge occurrence. The source
family's occurrence-selection discipline must be preserved separately.
Also, the paper's printed remainder-insensitivity condition only inspects
equation left-hand sides. In its literal finite profile, Transform may leave
a matching non-equation entry unselected, changing result counts. Complete
contextual coverage is a stronger, separately stated condition. These facts
prevent inferring a full-language specialization from similarly named rules.

Operational equations require a separate distinction. A source instruction
$a\longrightarrow b$ can execute while the structural constraint $a=b$ is
inconsistent for distinct constants. The core performs that change by
consuming and producing occurrences; it need not assert $a=b$. Likewise,
substituting only a selected replacement need not equal substituting its
entire surrounding context. The two readings agree when the context is
fixed by the substitution, and a checked example separates them otherwise.

\section{The atomic transform rule}

Let $O$ be physical occurrence identities, $T$ a declared term syntax and
$\sigma$ a substitution with a declared equality relation $E$. A state is
$(C,W)$, where $C$ is an equation store and $W\subseteq O$ is finite.
An occurrence has a syntax label $\ell(o)$. Its value is interpreted under
the same $\sigma$ as all other labels and equations.

A rule instance has lists of take patterns $t$, read patterns $r$, tells
$c$, checks $g$ and output terms $p$. A selection has lists of occurrences
$T_0,R_0,P_0$, in port order. Write
\[
 C' = C \mathbin{+\!+} \operatorname{eq}(t,\ell(T_0))
        \mathbin{+\!+} \operatorname{eq}(r,\ell(R_0))
        \mathbin{+\!+} c .
\]
The transition is
\[
 (C,W)\longrightarrow
 \bigl(C',\ (W\setminus\operatorname{set}(T_0))\cup
                         \operatorname{set}(P_0)\bigr)
\]
under the following obligations:
\begin{enumerate}
\item Port lengths agree, and selected output labels equal $p$.
\item Takes and outputs are individually distinct. Takes and reads are
present, no taken occurrence is also read, and outputs are fresh in $W$.
Reads may repeat an occurrence.
\item $\Sol(C')\ne\varnothing$.
\item Every $\sigma\in\Sol(C')$ satisfies every check in the original
physical world $W$.
\end{enumerate}

This is the existing \code{OccurrenceStep.Selection.Valid} judgment and
\code{Selection.result}, not a second transition relation. The rule is
already freshened at this level. Contextual activation gives each schematic
use fresh declarations while retaining explicitly shared caller captures.
Its endpoints also carry the scope inventory and actual solver return.
Space permissions are additional formation obligations; physical presence
does not supply authority.

For contextual terms, solvability must range over assignments into future
declaration families. In a constructor-free signature with only a binary
hole, $M(x,y)=M(y,x)$ has no endomorphic solution in its original family:
such a solution would require a closed term there. A fresh nullary helper
$H$ supplies the symbolic solution $M(x,y)=H$. Adding an unused nullary
declaration can therefore change fixed-family consistency while leaving
the original semantic problem unchanged.

The contextual judgment \code{ContextualGuardedStep.Valid} combines port
and resource conditions, the original all-future guards, and future-family
solvability. Its \code{Requirements} record omits solvability so that the
solver can still refute a candidate; it is not full admission. Positive
activation and information admission use the same future-family criterion.
Injective arity-preserving declaration extension preserves and reflects
admission, including output shape; no arbitrary substitution is given that
reflection law.

The one-frame monoid comparison remains a restricted profile. Ground terms
for every sort suffice to turn any future model into an endomorphic model.
The theorem \code{solvable\_iff\_consistent} in
\code{ContextualFrameUnifier} constructs that conversion.
This grounding condition belongs to that comparison, not to
general contextual admission. The helper-requiring consuming execution
above is admitted without it, while its original equation store still has
no endomorphic model.

\begin{theorem}[Contextual realization]
Within the declared structural-pattern fragment, semantic admission holds
exactly when the existing worklist has a completed admitted return at some
finite fuel. With the original scope coverage, the return constructs an
actual firing and its complete computed endpoint. Every completed firing
recovers semantic admission, including outside that complete fragment.
\end{theorem}

The worklist computes unification. Guard admission is a separate proved
obligation retained in the return; the theorem does not claim that
unification alone decides every possible guard.

The consuming control starts with occurrences $\{7,8\}$, uses the
nontrivial swapped-hole equation to consume $7$, and produces $99$.
Its computed target has occurrences $\{8,99\}$ even though the original
family has no endomorphic model. The fresh helper is an unconstrained
symbolic parameter, not a proved ground value or a discarded residual
equation. Additional solver declarations preserve solvability and all
equality, absence and count asks written over the original post-activation
interface. This does not grant invariance to observations inspecting the
private helper representation itself.

The proposed Formation--Resources--Information organization is dependent:
formation determines ports and scopes, physical selection determines the
equations, and those equations determine what the guards establish.
Splitting the predicate does not prove independence of three parameters.

\section{What actually commutes}

\begin{theorem}[Guarded interaction diamond]
Suppose two selected events are each admitted at the same initial state.
Suppose their physical commits are independent and their combined equation
problem is consistent. For each absence or count check, suppose every
occurrence taken or made by the other event is non-unifiable with its
pattern under the combined equations. Both event orders are admitted and
end with the same physical occurrences and the same equation solutions.
Every fixed continuation schedule therefore has the same admission result.
\end{theorem}

The proof transfers resource admission by the existing occurrence
independence theorem. Each later equation problem still satisfies the
earlier event's equations. Excluded matching writes preserve its queried
occurrences. The final equation lists differ only by order. The Lean entry
is \code{ConstraintMultisetRewriting.Core}; its substantive witness consumes
two distinct requests imposing $X=f(Y)$ and $Y=1$, shares a persistent read
(including repeated read ports), and retains equality and absence checks.
The final solution fixes $X=f(1)$ and $Y=1$.

This theorem compares endpoints and subsequent fixed schedules. It does not
identify intermediate effect histories or permit arbitrary atomic fusion.
Nor does it say that all useful interactions commute: one event can enable
another by adding information.

For contextual terms, a solver may replace its declaration family and
introduce private helpers. A principal answer $a$ to $C$ satisfies
\[
 \Models(\delta,C)\quad\Longleftrightarrow\quad
 \exists\rho.\;\delta=\rho\circ a
\]
on every original declaration, with $\delta$ ranging over every future
declaration family. Solving $C$ and then the residual of $D$ gives a
principal answer to $C\mathbin{+\!+}D$. Principal solutions produced in
the opposite order factor through each other on all original terms.
An isomorphism between their complete helper families is not required.

The frame-unifier theorem \code{Principal.solve\_orders} proves this comparison.
The update theorem \code{Update.completed\_orders} applies it to actual
completed worklist updates. Controls solve a swapped-parameter equation
$M(x,y)=M(y,x)$ followed by a body constraint, and the reverse order.
One construction changes a two-parameter declaration to a nullary helper;
the original term instances still agree. The executable worklist also
solves the constraint with body $\lambda z.z$, preserving its local binder.
Unfinished residual queues are excluded from these completion claims.

\begin{theorem}[Two actual contextual orders]
Let two original events be semantically admitted at the same endpoint.
Suppose each resource selection remains available after the other event,
the two joint contextual problems are solvable structural-pattern problems,
and each event's guards hold at its respective second admission cut.
Then both actual two-event paths exist. After exchanging the two private
event interfaces and fixing the original caller interface, every final
all-future equality, absence and occurrence-count question agrees.
\end{theorem}

The construction supplies both completed worklist returns, fresh scope
relocations, complete inventories and original firing edges. It does not
take either path as a premise. If $C_i$ is event $i$'s added equation list,
the independently prepared common problem is
\[
  \iota(C)\mathbin{+\!+}\iota_1(C_1)\mathbin{+\!+}\iota_2(C_2).
\]
The actual second problem is an injective declaration extension of this
problem. A completed first solver has an empty residual worklist, but its
original equations remain in the endpoint; they must not be erased.
The full final row inventory is likewise transported from the common frame.
Exchanging event-private summands exchanges the two contributions, fixes
caller captures and preserves the joint solution set. No isomorphism of
the complete solver-helper families is assumed.

The construction is \code{ContextualPairFrames.realize\_two\_orders};
its observation theorem reuses the existing contextual firing relation.
A concrete instance consumes different requests, uses two independent
views of the same schema occurrence in one event, and shares that occurrence
with the other event. The events jointly impose
\[
 M(x,y)=M(y,x),\quad X=\operatorname{pair}(Y,Y),\quad
 Y=a,\quad M(x,y)=\lambda z.z.
\]
Both paths exist, retain the incoming equation $Z=a$, and entail
$X=\operatorname{pair}(a,a)$ while preserving the lambda's own binder.
A separating model shows the second event alone cannot justify this value
of $X$. Thus retaining only its equation contribution loses an observable
constraint, even though both unifier worklists may finish.

The guard clauses remain semantic obligations at their original cuts;
resource independence alone does not prove them. The result concerns the
declared final questions, not observers inspecting solver internals or
intermediate effect histories.

\begin{theorem}[Causal contextual realization]
Suppose the first event is semantically admitted in the structural-pattern
fragment. Prepare each second-event input from its origin: a produced row
uses the first event's shared output; an older row uses its declared schema or
live view at the consumer port. If resources are available after production,
the joint problem is a solvable structural-pattern problem, output formation
holds, and the consumer's guards hold at its admission cut, then an actual
two-event path exists. Every final contextual Check agrees with the common
producer-origin presentation.
\end{theorem}

The theorem \code{ContextualPairFrames.realize\_causal\_pair} supplies both
solver returns and constructs the existing firing edges. The independent
input theorem is its specialization: when no consumer input is produced by
the first event, the two common equation problems are equal. Neither
construction introduces a second transition relation.

The causal control exports a private producer variable as live occurrence
$99$. The consumer is unavailable beforehand. After production, it either
reads $99$ twice or takes it, constraining its value to $\lambda z.z$.
Its guard counts one physical matching occurrence, even with two read ports.
The guarded take succeeds although the lambda is absent afterward; checking
the same count in the final world would reject this valid execution.
Two incompatible reads, demanding both the lambda and a distinct constant
from the one shared value, have no joint solution. Refusing that later
constraint problem does not undo the producer's valid publication: an
executed prefix and an atomic whole-chain refusal have different occurrence
worlds. General surrounding continuations, joined-schema guards, authority
integration and conditional atomic fusion remain separate obligations.

\section{Absence is a justified observation}

For a query $q$, the fixed-world absence guard is
\[
 \forall\sigma\in\Sol(C).\;\forall o\in W.\;
     \neg E(\sigma q,\sigma\ell(o)).
\]
Equivalently, each $C\mathbin{+\!+}[(q,\ell(o))]$ is inconsistent.
The transform separately requires consistency of its admitted equation
store, so vacuous entailment from an inconsistent store cannot admit it.

If $V\subseteq W$ is a scanned subset, a coverage certificate requires
every occurrence that matches in any solution of $C$ to lie in $V$.
Equivalently, every omitted occurrence is non-unifiable under $C$.
Under this condition, scanning $V$ preserves both absence and actual
occurrence counts. The existing complete first-order unifier constructs
such a candidate subset; an unfinished scan supplies no such certificate.
Full witness coverage is stronger than preserving absence alone: retaining
one of two matching occurrences preserves whether a witness exists, but
changes both the count and the retained provenance. These are different
observation contracts, even when they use the same matching relation.

The control Space contains $\operatorname{cancel}(0)$ and an open shared
variable $u$. Initially, absence of $\operatorname{cancel}(u)$ is unresolved.
Adding $u=1$ enables the selected guarded interaction; adding $u=0$ refuses
it. Inserting $\operatorname{cancel}(1)$ invalidates the first negative
observation, whereas an incompatible insertion does not. An empty partial
scan would incorrectly accept the guard and fails the coverage contract.

The Hyperseed d-calculus companion's~\cite{hyperseed} distinction between exhaustive
candidate coverage and an empty observation list is directly applicable
here. Existing Hyperseed visibility theory also characterizes when
restricted unary inference preserves all visible conclusions at a fixed
state. Joined rules and changing states require additional results.
Distinction-calculus evidence~\cite{dcalc} retains support and opposition separately;
missing evidence, explicit opposition, logical refutation, covered snapshot
absence and adding $u\ne0$ must retain their distinct meanings.

Adding a disequality is an information-changing operation. Equality of
successful final solution sets does not make irreversible publications
commute: a published output followed by inconsistency has a different
observable prefix from inconsistency before publication.

\section{Communication, quotation and the declared source profile}

The intrinsic rho binding-signature profile has a finite-path comparison
\code{source\_path\_iff\_firings} in the communication adequacy owner. Its
process classes use declared equations, including parallel commutative-monoid
laws and name quote/drop. Each source root compiles to an ordinary occurrence
rule; duplicate components retain separate identities. Parallel contexts
and ordinary capture-avoiding substitution are preserved. Successive firings
retain their incoming equations rather than reconstructing an empty store.
This is a process-class reachability result, not an executable complete
unifier for all source equations or an equivalence in arbitrary target contexts.

The source profile must be named precisely. The published rho presentation
in~\cite{lybech} seals quoted bodies against enclosing-name substitution.
Communication performs semantic substitution, which activates matching
received drops; a free Drop has no independent reduction there. The
intrinsic binding presentation instead substitutes through its ordinary
constructor arguments and can select a separate free-Drop rule. A checked
intrinsic receiver that refers to its input variable beneath quotation fails
the canonical rho quotation-scope check. These two carriers cannot be
identified by the unqualified constructor encoding.

A structural receiver interpretation now separates the received name's
ordinary name view from its active process view. It selects the view from
each occurrence in the authored body instead of asking higher-order matching
to infer a function from the compound spine $K[x,\operatorname{drop}(x)]$.
The interpretation commutes with contextual metavariable assignment,
including beneath nested input binders. The existing source quotation-safety
judgment gives a stronger boundary: compiled sealed code is independent of
the surrounding receiver environment. A nested-binder control agrees with
the independent semantic substitution; a literal quoted Drop stays literal,
and an outer-binder crossing beneath quotation fails the scope check.
These results are in \code{RhoContinuation},
\code{RhoContinuationSealing} and their controls. General agreement with
canonical source substitution, descent through source equations and actual
ordinary-rule execution remain required before claiming the full semantic
communication comparison.

An actual two-stage occurrence execution first publishes
$\operatorname{drop}(\operatorname{quote}(P))$ and then publishes an active
output $P$. A public absence query for $P$ succeeds at the intermediate state
and fails at the final state. Thus agreement after both stages does not
establish an atomic communication comparison in such observing contexts.
An encoding must justify private preparation, atomic publication, or an
explicit restriction on observations. This is an obstruction to that
unqualified staging claim, not to the transform mechanism itself.

The symmetric MeTTa-calculus manuscript has another independent boundary.
Its displayed REFL rule takes a derivation of the same reduction relation
as its premise, so reflection can be nested. The local executor instead
uses a COMM-only lookahead with fixed unification and context bounds.
\code{MeTTaCalculus.ReflectionClosure} states the displayed four-rule
closure with explicit unification and equivalence contracts. With the same
executable unifier, a two-level reflected COMM example has a derivation in
that closure while the current executor returns no result. This does not
supply complete alpha-equivalence or unifier semantics; it isolates why
that executor cannot serve as the unrestricted source specification.

Quotation remains a requirement of the Prime foundation. Existing candidate
owners define names for quoted terms and an explicit evaluation request.
The earlier guarded-rewriting reflection theory separates syntax and value
readers, but its abstract evaluation step takes an effects function as an
input. The mixed lexical-quotation construction below connects quoted lambda
code, its captures and explicit activation to retained transform execution.
Its quotation convention remains distinct from substitution-sealed rho names.

\section{Ordinary application through contextual transforms}

The independent source \code{LambdaContextualRung} has four operational
rules: beta contraction, reduction on either side of application, and
reduction beneath a lambda. Its object variables have intrinsic binding
positions. A one-hole context $K:\Gamma\Rightarrow\Delta$ records which
variables are available at the selected subterm, including binders crossed
on the way there. Such a context cannot erase or duplicate its hole.

For each selected context, the transform schema consumes
\[
 K[\operatorname{app}(\operatorname{lam}(B[x,\Delta]),A[\Delta])]
\]
and produces $K[B[A[\Delta],\Delta]]$. Here $B$ and $A$ are fresh private
contextual declarations; the ordinary body and argument are obtained by
unification with the input occurrence. No source evaluator or precomputed
reduct is an input to this rule. The produced term may contain compound
metavariable arguments; the input unification problem belongs to the proved
structural-pattern fragment for every ordinary source body and argument.

\begin{theorem}[Contextual lambda step comparison]
For every source context $\Gamma$ and source terms $s,t$, the four-rule
source reduces $s$ to $t$ exactly when a selected schema has a completed
actual contextual-worklist return from the encoded input occurrence and
its emitted value observes as the encoding of $t$. Every solution of the
admitted input equations gives that same reduct. Encoding commutes with
capture-avoiding object substitution and plugging a source context.
\end{theorem}

The theorem \code{LambdaComparison.source\_step\_iff\_executes} constructs
the forward solver return and recovers an ordinary source redex in the
reverse direction. The actual occurrence update consumes identity $0$ and
publishes identity $1$. The emitted row remains a live term in the computed
helper frame; interpreting its declarations gives the source reduct.
This solver theorem concerns one step over a represented source input.
Composition over retained inventories is supplied by a distinct semantic
result below; it does not assume that later raw problems remain structural
patterns merely because their represented values are ordinary lambda terms.

\begin{theorem}[Retained lambda execution]
The selected contextual-transform profile and the independent four-rule
source simulate each other step by step. From the encoded initial state,
finite paths have exactly the same represented lambda results. Intermediate
core states are the actual activation inventories: old equations, emitted
live terms, retained roots and the finite scope ledger persist. Every old
joint solution extends through each source contraction, preserving all old
term observations after the declared inclusion.
\end{theorem}

The theorem \code{LambdaExecution.finite\_execution\_iff} uses one local
declaration family indexed by the selected object's context and the
body/argument role. A later redex may therefore occur at a different binder
depth without changing the language of local declarations. Its transition
profile is generated by the context-indexed schemas above and the existing
contextual step; neither a source reduction premise nor a represented-output
premise appears in the transition definition. The proof recovers source
steps from admitted core steps and constructs core steps from source steps.

The control program first contracts beneath two binders, then beneath one.
The second transform consumes the first endpoint's emitted live term.
Both original matching equations remain afterward, and the occurrence
identity advances twice. This separates the construction from resetting to
an independently encoded source term between steps.

\begin{theorem}[Incremental solver realization of retained lambda paths]
The same finite source correspondence has an actual sequence of worklist
computations. Each completed prefix update extends across the fresh local
declarations without constraining them. Its computed answer transports the
new matching equations into the structural-pattern fragment. Continuing
the existing worklist completes at some finite fuel for every admitted
source contraction and retains a certificate for the entire original
equation prefix. Completed results recover source reductions without an
assumed valid source step; inert source variables are certified refusals.
\end{theorem}

The entry \code{LambdaExecution.finite\_execution\_with\_solvers\_iff}
records the actual calls in a solver history. The executable continuation
keeps any unfinished prefix residual alongside transported new equations;
fuel exhaustion never counts as completion. More generally, resumption
decides the original joint problem whenever this transported queue is a
structural pattern, even if its original syntax was not. For example,
solving $F(x)=x$ under a binder exposes a later
$F(\operatorname{pair}(X,a))=\operatorname{pair}(a,a)$ as a pattern problem
and computes $X=a$. Replacing its right-hand side by $a$ is refuted. Omitting
the earlier equation would admit the constant-function model instead.
The two-step binder-depth control also executes through the cached solver,
with both returned residuals empty. These are solver-termination results
for selected interactions, not normalization or fairness of lambda programs.

The rule family is indexed by explicit source one-hole contexts; a finite
installed evaluator, arbitrary enriched target contexts and resource-aware
evaluation strategies are additional constructions.

Controls distinguish $(\lambda x.\lambda y.x)z\longrightarrow\lambda y.z$
from the captured result $\lambda y.y$, and execute below lambdas and on
both sides of application. A source variable admits no invented step.
With $I=\lambda x.x$ and $\Omega=(\lambda x.xx)(\lambda x.xx)$,
$(\lambda x.I)\Omega$ can reduce to $I$, but the same unrestricted source
also permits repeatedly reducing $\Omega$. Consequently this comparison
does not supply Prime's demanded evaluation policy, correlated sharing,
effects or suspension. The demanded-machine comparison below treats those
behaviors separately; full beta reduction must not be substituted for
call-by-need semantics.

\section{Mixed quotation and activation with retained solver histories}

The independent \code{QuotedLambda} source adds
$\mathsf{quote}(t)$ and $\mathsf{activate}(t)$ to the intrinsic lambda
signature. Its active contexts close under application, lambda and the
argument of activation, but never enter a quotation body. The new root is
\[
 \mathsf{activate}(\mathsf{quote}(t))\longrightarrow t.
\]
Object substitution still traverses quotation to preserve its lexical
captures. A quotation can therefore be inert under reduction while changing
when an enclosing lambda supplies one of its captured variables. This is
not the substitution-sealed convention of the separate paper-faithful rho
profile. A concrete substitution distinguishes those contracts.

For each active context, beta and activation are ordinary transform schemas.
The activation hole has the entire ambient object context as its declared
arity. A nullary code hole would incorrectly require captured code to be
closed. Matching inspects rigid syntax; the rule output is a contextual
term, not the result of calling a source evaluator.

\begin{theorem}[Mixed finite execution]
Every finite reduction of the quoted-lambda source has an execution through
the existing contextual transform inventories with actual incremental
worklist computations. Conversely, every such selected execution from a
represented source term returns a source-reachable answer. All intermediate
equations, shared output terms and scope inventories are retained. Each old
joint solution extends through a source step without changing its old
observations.
\end{theorem}
The theorem is \code{QuotedLambdaRetained.finite\_execution\_with\_solvers\_iff}.
The independent source embedding preserves and reflects all finite pure
lambda paths, and the new represented-answer comparison agrees with the
previous \code{LambdaExecution} solver-path theorem. This compares answers;
it does not identify private declaration tables from the two clients.

A concrete mixed program evaluates
\[
 \mathsf{activate}((\lambda x.\mathsf{quote}(\lambda y.x))\,z)
 \longrightarrow \mathsf{activate}(\mathsf{quote}(\lambda y.z))
 \longrightarrow \lambda y.z.
\]
Its exact second endpoint retains both equations, and its solver history
continues the first completed cache. The result cannot be read as the
captured term $\lambda y.y$. An executable Lean probe completes both actual
solver calls. It also refutes attempted execution of held quoted code and
leaves a zero-fuel computation unresolved. The latter is not a negative
answer. A quoted looping term has no execution in this profile, and one
activation removes exactly one quotation layer.

These results close the lexical quoted-lambda comparison, not Prime's
complete reflection or evaluation policy. They do not provide a finite
installed context selector, demanded sharing, effects, unrestricted syntax
inspection, or composition with communication, complete querying and
authority. In particular, beta reduction does not identify quoted syntax:
the source has a beta step whose input and output have distinct quotations.

\section{An exposed instruction target for demanded execution}

The independently owned \code{CallTimeChoiceSpec} gives natural semantics
for computations, memoized and ordinary suspensions, forcing, constructor
cases, ordered equation dispatch, choice, failure, events, quoted code and
activation. Results retain their own branch stores. Its existing source
translation compares eager, call-by-name and call-by-need argument profiles;
none is silently selected as Prime's default here.

\code{CallTimeChoiceInstructions} exposes the recursive premises of those
judgments as finite root instruction families and explicit continuation
frames. An alternative frame retains the original store; sequencing forwards
each result's own refined store. Result concatenation preserves order and
duplicates. Neither instruction admission nor the executable next-step
function calls the source interpreter or assumes a completed evaluation.

\begin{theorem}[Complete source observations and explicit instructions]
A source judgment derives an ordered result list exactly when the
instruction execution reaches that list with an empty continuation stack.
The correspondence preserves the complete resulting stores and their event
traces. Every instruction agrees with the independently implemented
next-step function, so all finite instruction paths are executable.
\end{theorem}
The \code{CallTimeChoiceInstructions} owner contains:
\begin{itemize}
\item \code{source\_iff\_instruction\_path};
\item \code{instruction\_iff\_next};
\item \code{source\_iff\_executable}.
\end{itemize}
Missing references and exhausted fuel
retain unfinished control rather than producing a completed empty list.

Bookkeeping cannot diverge on its own. Between authored computation/force
boundaries, a lexicographic rank decreases: first the number of frames that
can launch further work, then the pending result-list traversal and remaining
stack. This is \code{administrative\_wellFounded}; it is not a normalization
claim for recursively defined source programs.

Executed controls distinguish two correlated shared-choice results from
four independent results, one shared effect from two repeated effects,
duplicate equation occurrences in their original order, unused recursive
work, inert quotation, repeated code activation, and residual noncode
activation. The source observation laws and code/heap names remain their
existing owners.

This instruction target is consumed by the per-cell construction below.
Structural environment lookup, cell operations and finite rule selection
are explicit data operations. Its comparison must preserve every inherited
cell at a branch, the complete return stack, and the ordered observation of
alternatives. Fresh branch names alone do not create inherited cells, and
a single-cell cache projection does not establish a whole-machine
simulation. No ordinary rule calls the instruction executor.

\subsection{Per-cell snapshot copying}

\code{SnapshotCopy} realizes copied stores by three ordinary rule families.
Preparation consumes a request, reads source readiness and metadata, checks
that the destination is unallocated, and produces a reservation and cursor.
Each copying event consumes the current cursor, reads one source cell, and
produces one destination cell and the next cursor. Completion consumes the
exhausted cursor and produces destination readiness and a reply. The payload
is closed data; this storage construction does not claim a new unification
fragment or call a whole-store copying function.

\begin{theorem}[Complete copying and all-execution progress]
Given an exact immutable source snapshot of $n$ cells, a fresh destination,
a present request and consistent equations, an ordinary execution copies
the snapshot in $n+2$ events. Every admitted execution from a valid prefix
with $r$ cells left has at most $r+1$ further events. If its destination is
ready, those $r+1$ events have occurred and the complete destination agrees
with the source, including cell indices, payloads and trace. There is no
infinite copying execution from such a prefix.
\end{theorem}
The central proofs are:
\begin{itemize}
\item \code{snapshot\_can\_be\_copied};
\item \code{Execution.ready\_iff\_completed};
\item \code{no\_infinite\_copy}.
\end{itemize}
For every source store, \code{NeedSnapshots} establishes these premises.
Each suspension retains its code, captured environment and memo state.
Explicit selected executions copy two
equal cached cells at distinct indices and an undemanded cell; all original
physical cells and intermediate equations survive. Wrong-index reads,
replayed cursors, incomplete completion and destination reuse are refused.

The invariant and observation boundary are essential. A context permitted
to read intermediate cells can see a partial copy. An external writer can
forge an exhausted cursor and thereby obtain readiness without coverage;
both separating examples are checked against actual ordinary firings.
Consequently this construction is not an atomic public snapshot operation
in arbitrary contexts. The joined instruction interpretation must protect
protocol state, establish freshness and preserve the chosen observations.
The global invariant below supplies these conditions for the generated
instruction catalogue. The linear number of copying events is a reference
cost, not an efficiency claim for a future persistent implementation.

\subsection{Store operations on the retained copy}

The copy proof now also derives unique destination metadata and reply
occurrences. The source snapshot need not satisfy these stronger conditions:
they follow from freshness and the actual produced rows. This matters when
an operation consumes a header; a duplicate old header must not remain.

\code{NeedStoreOperations} gives ordinary resource rules for allocation,
updating one memo and appending an event. Their parameters contain local
rows, code and environments, not a whole store or evaluator. A selected
linear command and the affected rows are consumed atomically. Other cells
remain in the frame. The independently defined source operations are
\code{Store.alloc}, \code{Store.setMemo} and \code{Store.emit}.

\begin{theorem}[Exact primitive store operations and retained frames]
Every source store operation has an ordinary firing from its editable
representation. Conversely, admission recovers its local source premises,
and every admitted result represents exactly the source operation's store.
All other snapshot identities retain their complete contents and physical
cell occurrences. Copying and an operation compose through the actual
intermediate world, preserving the original source snapshot.
\end{theorem}
The central results are:
\begin{itemize}
\item \code{admission\_iff\_matches} and \code{preserves};
\item \code{preserves\_other} and \code{copy\_and\_operate}.
\end{itemize}
The source's out-of-bounds memo update is a no-op. Its rule reads the exact
represented count and requires that the index lie outside it; it does not
infer absence from an incomplete scan.

A concrete retained execution changes one of two equal cached cells,
appends a repeated event, allocates an undemanded captured suspension, and
performs an out-of-bounds update. The other cached cell and original source
snapshot remain unchanged. Stale commands, old memo contents, stale counts
and a false out-of-bounds case are rejected. All prior equations survive.

These results concern a selected store operation. An evaluator must join
it atomically to the corresponding instruction control, preserving saved
continuations and result-associated snapshots. Installing the resource
families indiscriminately would not prove that comparison. The completed
command is not itself a grant of authority, and the copying rank does not
establish termination for arbitrary sequences of new source operations.

\subsection{Explicit controls and a physical execution episode}

\code{NeedControl} replaces store-valued positions in goals, continuations
and result lists by snapshot references. Each of the 42 source instruction
shapes has a plan with an input control, explicit data reads and either a
direct successor or a requested store operation. An edit continuation is
finite data: one marked position denotes the new snapshot, while marked old
references retain saved alternatives and other result stores. No callback
or whole source store is hidden in the continuation.

\begin{theorem}[Complete local control comparison]
A source instruction from an interpreted control exists exactly when a
supported plan has that input and its interpreted action has the specified
output. This covers the complete existing instruction grammar, including
missing-cell suspension. Every declared plan with its read facts yields a
source instruction.
\end{theorem}
The theorem is \code{source\_instruction\_iff\_supported\_plan}.
Its store interpretation is mathematical comparison data, not a premise
calling the source evaluator during rule admission.

\code{NeedMachine} puts heap and control rows in one ordinary occurrence
configuration. A direct start replaces the control and retains the heap;
an edit start produces a copy request and waiting control. Copy events read
that waiting control. Final commit consumes it with the actual requested
store resources, resumes the recorded continuation and does not reissue
the completed command. The operation and continuation are matched as part
of the control row.

A kernel-checked example performs thirteen actual firings. It copies the
original snapshot, changes the new memo from 7 to 42, resumes execution,
and returns the ordered pair of results 42 from the new snapshot and 7
from the saved original. The final physical heap is checked exactly and all
29 equations remain. The independent source instructions reach the same
two stores and values. Different operations, substituted continuations,
premature commits and replay are refused.

This selected physical execution starts at an internal commit continuation.
It is not yet a complete source-program entry comparison. The next results
provide the general projection, support and control laws needed to extend
it beyond a selected example.

\subsection{General projection, saved references and controlled completion}

The heap projection in \code{NeedMachineProjection} filters the actual
physical occurrences and accumulated equations of a global configuration.
It neither rebuilds a store from its values nor resets its constraint history.
Its admission theorem is restricted explicitly to the ground machine's
positive and heap-absence rules; it does not discard arbitrary guards.

\begin{theorem}[Controlled heap correspondence and edit completion]
Every controlled copy firing projects to an ordinary heap firing with the
same physical heap selection and endpoint. Conversely, every heap-copy
firing lifts in the presence of its waiting control and consistent global
equations. Complete heap-copy executions lift through their actual
intermediate worlds. From an exact source snapshot, fresh destination and
matching request and waiting control, a finite ordinary execution completes
the requested edit and resumes its recorded continuation. Every other
snapshot retains its complete contents.
\end{theorem}
The final commit consumes the completed reply without renewing it.
\code{published\_preserves} proves the edited snapshot law for this
publication rule by relating it to the existing renewable-command
operation. The extra command used in that proof is not a target execution
step or an equivalence for observers that count commands.
The principal results are \code{copy\_projects}, \code{copy\_lifts},
\code{copy\_execution\_lifts} and \code{edit\_block\_completes}.
The accompanying controls supply an inhabited finite physical instance
for every independently defined source store.

\code{NeedControlReferences} computes the finite snapshot support of
goals, result lists and every continuation frame. Decoding depends only on
this support. For all 42 plans, old output references and data-read
references come from the input control. If those input references lie below
the allocation frontier, filling the new-store marker and extending the
store interpretation preserves every saved old reference. Repeated
references are permitted: two branches may intentionally share a snapshot.
Cell indices inside values remain relative to their store; they are not
global snapshot names.

The fresh-resumption control returns 42 from the edited snapshot and 7
from its saved original. Reusing the old snapshot name instead returns 42
twice. A separate control retains snapshot 3 in a later continuation while
reading only snapshot 0 now. These examples show why freshness cannot be
checked against the current read set alone.

\begin{theorem}[Controller conservation and prepared waiting progress]
Every rule in the generated catalogue preserves the number of physical
controller occurrences. With exactly one controller, any firing from a
running state is a matching instruction start; any firing from a waiting
state is its matching copy or commit. From a valid prepared copying prefix,
there is no infinite catalogue execution retaining that waiting control.
\end{theorem}
The proof uses the existing copying rank after preparation and the actual
controller invariant. Two equal-labeled controller occurrences still count
as two; equality of labels cannot certify uniqueness. A concrete prepared
memo-update instance satisfies the no-infinite-waiting theorem.

The prepared waiting theorem alone does not bound an arbitrary sequence of
newly initiated edits. The following global invariant and rank supply the
stronger fixed-program comparison. Permissioned interference, open
constraints and changing definitions need their own joined comparison.

\subsection{Whole-program correspondence for the fixed-program client}

\code{NeedHeapFrames} distinguishes stored data below the allocation
frontier from one physical auxiliary occurrence. A running state has no
auxiliary occurrence. A waiting edit has exactly its request, current cursor
or completed reply, as appropriate. \code{NeedCopyInvariant} carries this
frame through the existing copying proofs, starting with the original
request. Publication consumes the reply and leaves only bounded stored
data. Equal labels cannot hide another reusable protocol token.

\code{NeedMachineInvariant} interprets only a finite allocated prefix of
the mathematical store map. Every live control, saved alternative and
result reference lies below that frontier, and every store in the prefix
has its actual snapshot representation. Running controls and waiting plans
remain unique. An instruction start either replaces the running control
or starts a waiting edit; copying preserves the saved stores; final commit
publishes the exact edited store and resumes the recorded continuation.
These properties hold for every admitted firing of the generated catalogue.
The independent source judgment is \code{CallTimeChoice.Derives}.

\begin{theorem}[Two-way finite execution from ordinary entry]
For every fixed source program, initial store, environment and authored
computation, the source judgment returns an ordered list of results exactly when an ordinary
catalogue execution from its physical entry world reaches a represented
running endpoint decoding to that list. The comparison preserves the full
store associated with each result, including suspensions, captured
environments, memo contents and ordered effect traces.
\end{theorem}
The principal results are \code{Running.source\_path\_iff} and
\code{NeedMachine.source\_derivation\_iff}. The latter uses the existing
source-to-instruction theorem and the new physical path comparison; the
target observation is an ordinary path with represented stores, not a
definition calling the source evaluator. Entry contains the source data
and one running control, with no preinstalled copy request.

Every physical step either realizes a source instruction or preserves the
source control and strictly decreases remaining copy/publication work.
Interpretation and remaining work are independent of the choice of
proof-side model witness: metadata and cells determine each referenced
store, while the controller and physical protocol token determine its
phase.

\begin{theorem}[Composed administrative progress]
Between source-work boundaries, a lexicographic rank decreases: the
existing source continuation rank first, then remaining copy/publication
work. Every copying transition is included even when its decoded
continuation is poised to perform source work. Consequently the generated
machine has no infinite execution consisting solely of this bookkeeping.
An unfinished copy has an available next firing; a missing source cell
suspends the corresponding running catalogue state.
\end{theorem}
This is \code{NeedMachine.administrative\_wellFounded}, not a termination
or fairness theorem for arbitrary recursive source programs.

Concrete ordinary-entry observations are exact: shared choice returns two
diagonal pairs; independent suspensions and ordinary delay return four;
a shared effect occurs once and repeated delay emits it twice. Duplicate
alternatives and definition occurrences retain order and multiplicity.
Unused recursive work stays unforced, quotation stays inert, repeated
activation repeats its declared effects, and noncode activation yields
the declared residual. The comparison also excludes every different
complete observation, using determinism of the independently specified
ordered-result judgment.

This closes the per-cell realization of this fixed-program, ground client.
It does not close full Prime: its combination with contextual open
unification, the quoted-lambda binding profile, changing definitions,
queries and external authority still requires explicit comparison. Nor
does the reference's retained snapshot prefix prescribe physical storage.
Representation refinements must preserve their declared source observations;
raw private heap and allocation-history observations are different contracts.

\subsection{A concrete connection to native behavioral logic}

OSLF derives logical observations from a declared operational
presentation~\cite{logic-distributive}. A source language description and an
ordinary transform execution are therefore compared by a semantic map;
the availability of both descriptions does not make them isomorphic.

\code{NeedCopyObservation} connects the controlled-copy subsystem directly
to the existing native modal semantics. On one side, states are actual
global configurations with consistent equations and the matching waiting
control. On the other, they are the actual projected heap configurations.
Events on both sides are ordinary copy firings. The state map forgets the
controller, and its event map uses \code{copy\_projects}.
\code{projection\_source\_lifts} uses \code{copy\_lifts} to lift every
outgoing heap event at the observed state.

\begin{theorem}[Exact future behavioral observations of copying]
Let $h$ be this heap projection. For every formula $\varphi$ using only
future modalities, with positive fixed-point bodies, and every environment
$\eta$ of heap predicates,
\[
  h(w)\models_{\rm heap,\eta}\varphi
  \quad\Longleftrightarrow\quad
  w\models_{\rm controlled,h^*\eta}\varphi.
\]
This includes nested least and greatest fixed points. The proof reuses the
native modal transport theorem with outgoing endpoint coverage derived
from actual ordinary firings.
\end{theorem}
The result is \code{future\_formula\_iff}. A concrete preparation,
cell-copy and completion execution witnesses
$\mu X.(\mathit{ready}\lor\Diamond_{\rm future}X)$ on the controlled side.
The observation restriction matters: this compares heap predicates and
future behavior of copying, not raw controller observations, spatial
decomposition, incoming history, event counts, fairness or whole programs.
It is a semantic theorem, not a newly installed runtime model checker.

The native framework distinguishes directions. Its adjoint OSLF box reads
incoming predecessors, while a universal future box reads successors.
At the concrete initial copy state, the former applied to false is true
because there is no incoming copy; the latter is false because preparation
can fire. Neither one-step operator alone means temporal ``always'' or
``eventually''. Fair eventual service additionally requires a scheduling
and progress contract.

The source-name reuse corollary is deliberately narrower than physical
implementation equivalence. If the requested source name does not occur
among a continuation's retained old references, overwriting that entry in
the mathematical store interpretation gives the intended decoded
continuation, equal to fresh placement. An independent reader of that old
name distinguishes the two placements, returning 42 after overwrite and 7
after fresh placement. A physical refinement must therefore account for
all permitted observers, aliasing, equations and effects.

The reference copying construction does not prescribe a storage backend,
an append-only history, or copying on every language operation. Direct
updates, trails, persistent representations and other implementations
remain candidates for explicit semantic and cost comparisons. A stored
snapshot is also only one component of a behavioral state: active rules,
control, permissions and definitions can affect its transitions.

\subsection{Behavior at source-instruction boundaries}

\code{NeedMachineEpisodes} defines one instruction episode entirely from
ordinary firings: a direct start, or an edit start followed by only its
copy rules and the matching commit. A source-instruction premise is absent
from this definition. Every episode expands to an actual catalogue path.
The global invariant proves that it realizes exactly one independent
source instruction; conversely, every source instruction from a represented
running state has such an episode. The original data admission cut and
the full intermediate equation histories are retained.

\code{NeedInstructionObservation} uses these episodes as its target
reduction span and the existing instruction relation as its source span.
Decoding represented running states defines the state map; the proved
episode comparison defines its event map and outgoing endpoint lifting.

\begin{theorem}[Exact future instruction observations]
For any fixed source program and represented running state, source
predicates and future-directed modal formulas with positive fixed-point
bodies have the same truth value in the instruction span and the actual
episode span. This includes nested least and greatest fixed points.
Native reachability agrees with a physical catalogue path to a represented
running endpoint satisfying the source predicate.
\end{theorem}
The central results in \code{NeedInstructionObservation} are:
\begin{itemize}
\item \code{projection\_source\_lifts};
\item \code{future\_formula\_iff};
\item \code{reachability\_iff\_physical\_path}.
\end{itemize}
Ordinary-entry shared choice
instantiates the native reachability formula for its two diagonal results.
The corresponding four-result resampling formula is refuted. A concrete
memo edit is one episode but five raw firings; no single raw firing reaches
its published endpoint.

\begin{theorem}[Incoming-box boundary]
Uniform preservation of the predecessor-universal box by this projection
is equivalent to incoming endpoint lifting. That additional property fails
for a concrete one-definition program: a fresh physical entry at the
function body has empty matching-equation history and no incoming episode,
while its decoded source control has a valid calling predecessor.
\end{theorem}
Thus the incoming box applied to false holds at that physical entry but
fails at its source image. The proofs are
\code{incoming\_box\_transport\_iff} and the accompanying
\code{incoming\_box\_is\_not\_preserved} control. This is an obstruction
to an unrestricted history-forgetting comparison under the stated
observations, not an obstruction to transform execution.

Logical source predecessors and physical execution history must therefore
have declared interpretations. A different semantic quotient or history
observation would need its own correspondence proof. No storage policy is
selected by these results. Spatial decomposition, event multiplicity,
scheduling fairness, external permissions and definition revisions remain
separate comparison obligations; no new executable runtime model checker
is supplied by the semantic transport theorem.

\section{Resolved clauses and subsequent admission}

\code{ContextualClauseSourceState} defines a source world of contextual
declarations, equations, stored schemas or shared references, and physical
occurrences. Its update is constructed directly from an admitted resolved
clause. Query variables are shared across that clause's ports; each use of
a stored schema has its own private variables and latent equations.
The source continuation retains its returned query environment.

\code{ContextualClauseStateRefinement} compares this independent update
with the actual clause compiler through an injective capture frame $h$.
The comparison preserves information, solvability, physical effects,
conditional row observations and the returned environment. Native helper
allocation may change declaration families. Unfinished solver equations
remain attached to reusable outputs.

Subsequent joins require more than separate existential matches: all ports
must use the same query assignment and caller valuation.
\code{ContextualClauseWitnesses} supplies these shared assignments together
with private witnesses for each port in one future family. The native/source
row comparison preserves these simultaneous realizations, including each
row's latent equations.

\begin{theorem}[Subsequent resolved-clause admission]
After an actual compiled clause, let $r$ be another resolved clause over
the independent source result, with a selected occurrence list $s$.
Transport $r$'s captured declarations through the actual frame $h$ while
retaining its query variables. The transported clause is admitted in the
native result exactly when $r$ is admitted in the source result.
The equivalence preserves selected occurrences, joint solvability,
shared equations, and equality, absence and count guards at the same
observation world. For every finite solver budget, source admission is
equivalent to existence of the corresponding next compiled firing at
that budget, retaining any residual equations.
\end{theorem}
The owner is \code{ContextualClauseSequencing}.
Its admission result is \code{next\_admitted\_iff}; the compiler result is
\code{next\_compilation\_iff}.
Unused private arities are supplied by actual captured declarations;
the proof assumes no inhabitant at an arbitrary sort.

A concrete first clause reads a stored function
$\lambda z.\mathsf{pair}(z,2)$, applies its contextual body to $1$,
uses the computed key to obtain answer $7$, consumes a permit and
publishes a reusable answer. A second clause uses the retained binding,
reads that published occurrence twice, checks the binding against $7$
and publishes at another space. Its actual compiled output matches
exactly answer $7$. Replacing the second read by a demand for $9$
admits neither a source firing nor a compiled firing.

\subsection{A comparison preserved by clause execution}

Further execution may change the private declaration families differently
on the two sides. The world comparison therefore fixes an exposed
environment while allowing different private representations. Every model
of either world has a corresponding model of the other in a fresh future.
The models agree on the exposed bindings and on conditional matches of
every present occurrence in all further fresh futures. Physical occurrence
identities remain equal.

Fresh witnesses of finitely many nonbinding row matches can be amalgamated
without refining their shared query assignment. This supplies one joint
model for all ports. It does not permit each port to choose its own query
bindings. Fresh declaration extensions are logical renamings; they do not
specify a runtime snapshot or copying operation.

\begin{theorem}[Preservation of the world comparison]
If two worlds satisfy this comparison, admission of the same authored
resolved clause through their exposed environments agrees. Their source
updates again satisfy the comparison, with the query bindings added to
the exposed environment. An actual compiled firing preserves the same
comparison between the independent source update and its native result.
\end{theorem}
Retained rows use the old comparison. Live outputs use equality of the
returned environment. Reusable outputs use the two-way correspondence
of complete conditional source instances, including their latent equations.
The source-update result is in \code{ContextualWorldClauseResults}.
The actual compiler-step result is in \code{ContextualWorldCompilerSteps}.

\subsection{Finite authored programs}

\code{ContextualClauseSequences} defines a finite resolved program fragment:
return a contextual value, or execute an authored clause and continue with
its query bindings and old exposed environment. Its source execution uses
the independently constructed world update. Its native execution uses
actual compiler firings. Both record an ordered list of physical selections.
This is mathematical syntax for the comparison, without a textual scope
or search-order policy.

\begin{theorem}[Finite clause-sequence adequacy]
For source and native worlds satisfying the world comparison, a finite
authored clause sequence has a source execution with selection trace $\tau$
and final conditional returned value $v$ exactly when it has a native
compiler history with the same sequence, trace and returned value.
Every native step carries an actual worklist computation; remaining
equations are retained. The statement preserves duplicate reads and
physical occurrence distinctions.
\end{theorem}
The theorem is \code{source\_iff\_native} in
\code{ContextualClauseSequenceComparison}. Its induction uses preservation
of the world comparison at each step and exposed-value instance agreement
at return. Neither native admission nor a rule calls a whole source evaluator.
Semantic admission evidence remains explicit; this does not provide a
total automatic guard decision procedure or a completing solver budget.

The computed two-clause example is now also an authored sequence with a
source run and an actual native history, returning $7$ with exactly the
declared physical selections. Its conflicting second read has neither
history at that trace. A separate control proves that replacing the
exposed answer $7$ by $9$ is not a world comparison even if the physical
world stays the same.

Composition of complete answer censuses, arbitrary active evaluator and
recursive/choice contexts, and compilation of family-changing private
effect episodes into one public update remain separate obligations.
The finite sequence performs one atomic firing per clause; it does not
fuse the entire sequence into one firing. Additional observations of
private representations need their own comparison. No textual scope
default or physical storage policy is selected here.

\subsection{Publication prefixes in an atomic computed body}

An ordinary schema published before a later caller binding retains its
earlier meaning. A private body that first publishes
$\mathsf{box}(X)$ while $X$ is free, then binds the caller's $X$ to $1$,
still has a reusable schema matching $\mathsf{box}(2)$. Publishing
$\mathsf{box}(X)$ against the final joint problem loses that answer.
Likewise, a check that $X$ was already $1$ before the binding cannot be
discharged by evaluating it only against the final information.

\begin{theorem}[Independent prefix isolation]
Let $P$ be the saved prefix problem and $Q$ a solvable later problem.
Embed their declaration families into disjoint summands, with inclusions
$\iota_P$ and $\iota_Q$, and form
\[
 R=\iota_P(P)\mathbin{+\!+}\iota_Q(Q).
\]
The conditional instances of $\iota_P(v)$ under $R$ are exactly the
instances of $v$ under $P$. An original-cut guard on rows transported
through $\iota_P$ is entailed by $R$ exactly when it was entailed by $P$,
at the same watched occurrence world. Arbitrary contextual arities and
latent equations are retained.
\end{theorem}
The owner is \code{ContextualPrefixIsolation}. A model of the saved
prefix extends by a separate model of $Q$; unused arities receive actual
fresh generators. This gives both instance preservation and reflection
of universal guard asks. Solvability of $Q$ is essential to prevent a
vacuous conclusion.

The native construction is in \code{ContextualPrefixNativePublication}.
Every finite budget has an actual worklist result when both factors are
solvable. Its reusable output retains the actual residual equations and
has exactly the saved prefix's instances. A concrete saved condition
$Y=\mathsf{box}(Z)$ still admits $\mathsf{box}(2)$ after the ambient
caller's own $Z$ becomes $1$, and excludes bare $2$. A zero-budget result
retains a nonempty three-equation problem, with that same exclusion.

For the open-schema example, a single existing compiler firing using a
separate private schema variable has the same exposed post-state
observations as the two private source steps. It preserves the caller's
binding as well as the earlier schema's complete instance set. The
construction selects no parent-space snapshot or physical copying policy.
The concrete encoding is supplemented by a general selected-body lowering.
\code{ContextualSourceSchedules} represents finite private clause schedules
without assuming admission. Compilation in
\code{ContextualScheduleCompilation} collects binding-equation deltas,
guards at their original cuts, and production origins retaining their own
declarations, conditions and values. It computes a net occurrence update
and an initial resource precondition rather than executing the source
schedule as a callback.

\begin{theorem}[Exact admission of a compiled private schedule]
The compiled shapes, resource precondition, joint information and original
guard cuts hold exactly when the selected source schedule is admitted.
For an admitted schedule, the compiled net resource update is enabled
and yields exactly the source's final occurrence set.
\end{theorem}
The results are \code{Schedule.compile\_exact} and
\code{Schedule.compile\_effect} in that compiler. The controls retain the
earlier schema's answer while fixing the terminal caller, and reject both
a conflicting second read and an earlier guard that only later becomes
true. The latter has consistent terminal equations and resources, so those
alone do not certify admission. Atomic assembly of the production origins
into an actual firing, followed by evaluator and complete-census
composition, remains open.

\subsection{Bound allocation names at the atomic boundary}

A source body can consume occurrence $0$ carrying $1$ and later reuse
the number $0$ for a row carrying $2$. Both private steps satisfy local
freshness. A single existing firing from the entry world cannot create
that changed row at $0$, since its output identifiers must be fresh at
entry. Its unchanged final occurrence set does not express the row
replacement. \code{ContextualAtomicIdentityControls} proves this
separation and constructs an actual solver-backed repair that consumes
$0$ and produces the new row at fresh identifier $1$.

The general resource comparison uses bound names for internal creations,
leaving the existing source and step definitions unchanged. Entry names
are embedded as $\mathsf{inl}(o)$; a creation at position $k$ receives
$\mathsf{inr}(k,o)$. The current map is injective and erases back to $o$.
It can therefore distinguish a reused number's new creation from its
old entry occurrence. One global permutation fixing the old number would
not make this distinction.

\begin{theorem}[Fresh allocation comparison]
Every admitted finite source schedule lifts to an executable resource
schedule with these fresh creation names. Erasure recovers its exact
ordered take, read and produce selections, and its final occurrence set.
The lifted schedule has one enabled net occurrence commit whose result
erases to that same source result.
\end{theorem}
\code{OccurrenceAllocationRenaming} proves the resource comparison;
\code{ContextualScheduleAllocation} applies it to the independently
admitted contextual schedules. The entry names and internal creations
belong to separate summands. These are comparison names for physical
occurrences, not changes to object-language numbers, space handles,
logical variables or a required runtime representation. Conditional rows
and original guard cuts still require their atomic assembly comparison;
the resource theorem alone does not supply that result.

\subsection{Conditional rows and guard cuts under allocation renaming}

\code{ContextualOccurrenceRenaming} transports conditional row regions
through an injective physical map. The represented region is exactly the
image of the original region; absence and cardinality checks therefore
agree. The row correspondence is needed only on live occurrences and
retains their latent equations. Identifying two equal-valued schemas
would preserve neither their physical multiplicity nor their count.

\begin{theorem}[Original cuts and renamed clause publication]
Every saved guard cut in a selected private body has the same all-future
meaning after renaming by the allocation map at that cut. For a single
clause, renaming retains its complete joint binding problem and source
update row meanings. Source admission agrees with renamed admission
together with the original source freshness condition. Every admitted
renamed clause has an actual native compiler firing at every finite
budget, whose conditional row observations agree with the original
source update at the corresponding current occurrence names.
\end{theorem}
The saved-cut proof is implemented in
\code{ContextualScheduleCutRenaming}.

The per-clause result is implemented in
\code{ContextualClauseOccurrenceRenaming}. Original freshness is
essential: a renamed output may be physically fresh even when the
source tried to overwrite an occupied identity without consuming it.
That invalid source operation remains refused.

The controls retain six pending equations from two uses of a conditional
schema, construct an actual renamed publication with the source's row
observations, and refuse a conflicting read. Earlier false guards stay
false after later bindings. Two schemas with equal values still count
as two occurrences; collapsing their names changes that count. These
results leave the general assembly of all surviving production origins
into one public firing open.

\subsection{Constructed row payloads for the complete private body}

The result data can now be assembled without reading the final source
world's rows.

\code{ContextualScheduleOrigins} traverses initial input
rows and authored productions. It records retained input with its capture
transport, ordinary schemas with their own saved conditions and values,
and live expressions. Later equations act on the explicit capture maps;
they do not enter an earlier ordinary schema's saved problem.

\code{ContextualConditionalRowPacking} puts these independently indexed
rows into one private bank. Both values and latent equations use the
same inclusion; shared captures stay shared. A model for unused private
arities is supplied by fresh formal generators in an extension, rather
than by arbitrary inhabitants. Equal-valued occurrences are not merged.

\begin{theorem}[Conditional observations of the assembled result]
For every admitted selected private body, the compiled origins and packed
rows have exactly the final source rows' conditional instances in every
later declaration family. The assembled information is exactly the
source's final information, and the renamed net resource update gives
the corresponding final occurrence set. At every finite solver budget
there is an actual worklist result over that information. Its constructed
row store preserves those row observations and returned-value instances,
including when the worklist retains pending equations.
\end{theorem}
The data assembly is implemented in \code{ContextualScheduleRowAssembly}.
The bounded computation is implemented separately in
\code{ContextualComputedRowAssembly}.
Even an unfinished admitted result has a solvable retained equation queue.
Unselected dormant row conditions are not conjoined with that queue.

Controls retain the earlier schema's answer $\mathsf{box}(2)$ while
fixing the caller to $1$, and keep one pending equation at zero fuel.
A body computing a key through stored code accepts answer $7$ and
excludes $9$. A separate capture control links a row-private variable
to a shared reference; packing preserves the link, whereas dropping
its latent equation invents an answer.

This constructs the complete result data for the selected ordinary/shared
clause profile. It does not identify that construction with an ordinary
single-clause native firing. Such a lowering must account for per-output
creation conditions and original guard contexts. No evaluator callback,
new firing-graph edge or physical copying policy is introduced here.

\subsection{Uniform publication of independently scoped birth problems}

The ordinary native publication form need not be changed merely because
outputs have different creation conditions. Each produced descriptor
retains a solvable creation problem. Put these genuinely private copies
in separately tagged declaration families, alongside the current caller
problem, and solve their conjunction with the existing bounded worklist.
The descriptor also retains a decidable name-equality dictionary for
that actual computation.

\begin{theorem}[Uniform private-condition representation]
For an admitted selected body, every current caller model extends to
the combined problem with its assignment fixed. Each creation value has
exactly its original conditional instances under the combined problem.
At every finite budget, one actual returned update can therefore supply
the unchanged ordinary publication constructor for every birth, with
those same instances. Independent satisfiable blocks also preserve a
guard ask at its own original watched world.
\end{theorem}
The general block comparison is in
\code{ContextualIndependentProblems}.
The body-specific construction is in
\code{ContextualUniformBirthProblem}.

One actual ordinary native firing publishes two schemas together: the
earlier schema admits $\mathsf{box}(2)$, while the later one admits
$\mathsf{box}(1)$ and excludes $\mathsf{box}(2)$. Both use the same
existing solver and publication operation, at every finite budget.

Private independence is essential. In a separate control, two row
conditions share one captured variable. Their conjunction is globally
satisfiable when that capture is $9$, but suppresses an otherwise valid
row at capture $7$. The private-block theorem does not apply to that
shared condition. Dormant input rows remain conditional, and explicitly
shared captures are preserved under the same caller model. The general
ordinary firing lowering and conditional captured-template integration
are still open.

The production lookup now connects the finite net output set to actual
creation descriptors. Every physical net output has a surviving source
production index. The resulting ordinary and shared output forms, together
with retained conditional inputs, preserve all source row instances under
the returned caller environment. A four-clause control recycles source
occurrence $0$ twice; its surviving value is $3$, with creation name $(3,0)$,
rather than either earlier value.

\begin{theorem}[Joint information preserves original guards]
For any selected body, place each original guard problem in a separate
private declaration block beside the current and creation problems.
The joint problem is solvable exactly when the terminal source problem
is solvable; no guard-success assumption is needed. Shapes, the compiled
resource precondition, joint solvability and all original-world guard
views are together equivalent to source admission. Every actual bounded
worklist return preserves and reflects these guard verdicts, even when
its remaining equation queue is nonempty.
\end{theorem}
The production connection is in
\code{ContextualProductionAllocation}.
The guard comparison is stated once for compiled bodies in
\code{ContextualNativeBodies}; \code{ContextualGuardedScheduleProblem}
presents selected ordinary/shared bodies as compiled bodies.
The controls include a consistent final problem with an earlier false
guard, and an admitted computed join with a true answer guard. Representing
the saved guard worlds in a native entry inventory remains a separate
operational obligation. These logical declaration blocks prescribe no
physical snapshot or copying policy.

\subsection{Entry-world regression and data dependencies}

Row comparison can precede admission: output shapes suffice for its
conditional instances, and resource executability suffices for its
physical projection. The compiler therefore constructs guard views even
for a candidate with a refused guard. At the same information cut,
conditional matching obeys the existing occurrence-query equation
\[
 R_{W\setminus T\,\cup P}
   = (R_W\setminus R_T)\cup R_P.
\]
When retained row meanings agree, absence checks both surviving inputs
and creations, while a fresh update has count
$|R_W|-|R_T|+|R_P|$. The original clause's own port and tell equations
remain in its guard problem. These constructions are in
\code{ContextualGuardEntryViews}.

A pointwise truth query at one marker does not validate the row data
captured by its predicate. A native control inserts a matching fact and
still satisfies the marker query, while its resulting world violates
the intended count. Native query labels now expose the actual conditional
row as well as its occurrence identity. A complete query of that row
region detects the matching insertion and permits an unrelated one.
The adapter reuses the existing predicate-query record; it is in
\code{ContextualPredicateQueryAdmission}.

\begin{theorem}[Principal count-query translation]
Let an actual completed update of $C$ return principal assignment
$\theta$, and let $S_\theta$ be the matching occurrence set under
$\theta$. Then
\[
 \operatorname{Ask}_{C}(\#p=n)
 \quad\Longleftrightarrow\quad
 |S_\theta|=n\ \land\
 \operatorname{Ask}_{C}(\operatorname{region}(p)=S_\theta).
\]
The query reads actual conditional row labels. Private matching witnesses
transport through further assignments, and every model of $C$ factors
through $\theta$. Thus $S_\theta$ is included in every model's matching
set; equal finite counts force equality of those occurrence sets.
\end{theorem}
This is proved in
\code{ContextualPrincipalCountQueries}.
Its native instance counts two identical conditional schemas at distinct
occurrences and publishes value $7$. A separate control has completed
unification and an empty principal table, yet a later caller assignment
creates a match; the complete query remains necessary. This translation
does not erase a pending equation queue or establish a total guard
decision procedure. A fixed matching set is specific to this
completed-principal profile; a conditional data table permits a more
general count comparison below.

\begin{theorem}[Native binding preserves every source model]
Let a finite selected private body have the stated port shapes and an
executable resource schedule. Its net take/read patterns and one explicit
conditional control row determine an ordinary native binding problem.
Every model of the terminal source information extends to a model of
that actual native problem through an injective fresh declaration frame.
The current source assignment is retained on every original name.
Conversely, every native binding model projects to a source model.
\end{theorem}
The construction is stated once for compiled bodies in
\code{ContextualNativeBindingProblem}.
Input witnesses are extracted from source ports in
\code{ContextualScheduleInputModels}; unused stored conditions are not
conjoined. Packed inputs retain their conditional instances, and the
initial declaration extension preserves the entry problem. The control
row is explicit program data. Erasing it is a declared projection, not
equality for observers of the complete inventory.

The concrete body opens a stored lambda, computes a lookup key, consumes
a permit and reuses the answer twice. Every native binding model retains
answer $7$, and incompatible reuse demanding $9$ has no model. Adding an
unused inconsistent schema leaves the binding problem consistent.
A separate control still has a binding model when an earlier guard is
false. Thus the binding comparison alone does not discharge guard admission.

\begin{theorem}[Original-cut native equality and absence]
Every authored cut has a finite entry descriptor with the same information
problem, checks and guard meaning. Its occurrence world is
$(W\setminus T_g)\cup P_g$, where $W$ is the entry world and $T_g,P_g$ are
the compiled prefix delta at that cut. In the actual native binding
problem, equality uses the original cut's copied terms. Absence is
equivalent to the conjunction of a complete query of surviving actual
entry labels and a finite absence check on the authored creations $P_g$.
\end{theorem}
The source descriptors are constructed in
\code{ContextualEntryGuardCompilation}.

\code{ContextualNativeGuardViews} resets actual row captures to the saved
cut, retaining private schema conditions. Every original-cut model extends
through the native information, so later equations do not erase its
countermodels. The actual-row predicate uses the supplied label data;
the creation check reads only compiled prefix descriptors.

The controls reject an earlier absence of $\operatorname{box}(2)$ from a
live $\operatorname{box}(X)$ even when a later clause sets $X=1$.
The same equation in the guarded clause itself correctly admits the check.
A matching insertion invalidates the query, while an unrelated insertion
does not. An empty entry inventory can still produce a matching schema
inside the private prefix: the entry query alone passes, and the creation
check restores refusal.

The comparisons fix the constructed entry world and compiled physical
selection. Production meanings and conditional captured outputs are joined
in the one-firing comparisons below; program-context observations remain
separate. No parent-space copying policy follows
from the logical declaration reset.

\begin{theorem}[Conditional count-table translation]
A finite table $T$ stores row data on support $D$. For each admissible
candidate $a$, its query validates equality of the actual and recorded
matching predicates on $D$, and rules out actual matches outside $D$.
Whenever this query succeeds, the actual matching region is exactly
\[
 M_T(a)=\{o\in D\mid\operatorname{matches}(a,T(o))\}.
\]
At an authored guard cut with prefix removals $T_g$ and creations $P_g$,
the source count guard is equivalent to that actual-data query together
with the finite check
\[
 |M_T(a)|+|M_{P_g}(a)|=n
\]
for every model of the actual native binding information. The table reads
only surviving entry occurrences; creations retain their own authored
conditional meanings. No completed principal assignment is required.
\end{theorem}
The reconstruction handles arbitrary supplied row labels and worlds.
Its generic owner is \code{SymbolicPredicateTable}.
The original-cut comparison is in
\code{ContextualSymbolicCountQueries}. A suitable support is the subset of
the finite inventory that matches in some admissible model. This is a
semantic construction, not an algorithm completing arbitrary higher-order
matching. The read validates matching behavior, not arbitrary payload or
weighted observations.

In a concrete contextual constraint, the unknown two-parameter term
context satisfies $F(0,1)=0$. Every solution is the constant $0$ or the
first projection, including in arbitrary future declaration families.
Consequently $F(1,0)$ matches exactly one of the stored literals $0,1$,
but its matching occurrence differs between these solutions. A fixed
successful-match table cannot work, and no completed principal update
exists. An actual firing at every finite worklist budget retains the
constraint and publishes $7$; its published row denotes that value.
The symbolic table validates at the same information cut. A changed
relevant payload or a new matching duplicate occurrence refuses its read.
A separate native compiler instance counts
$\operatorname{box}(1)$ after the guarded clause's own equation $X=1$,
accepting count one and refusing count zero.

These count comparisons retain the fixed constructed selection and the
stated source shape and resource hypotheses. Joining exact publication,
captured outputs and full program observations to the same actual firing
remains open.

\begin{theorem}[Native birth instances and conditional publication]
Under the selected body's shape and resource hypotheses, and consistency
of its terminal information, the actual native binding problem preserves
and reflects every original value instance in a saved birth block.
An actual conditional firing whose declared ordinary output is that
birth value publishes exactly its original conditional instances.
Shared outputs retain the current caller interpretation; retained input
metadata retains its original conditional instances under the returned
caller. The target information is satisfiable and projects to the source
terminal information.
\end{theorem}
The comparison in \code{ContextualNativeBirthPublication} first extends
each given compiled model through the actual selected input conditions.
Independent satisfiable birth blocks preserve the selected birth value.
It then applies the existing conditional producer's ordinary, live and
retained-row laws. Physical presence is determined separately by the
same firing's resource result.

A concrete native producer constructs an earlier reusable
$\operatorname{box}(X)$ followed by the caller equation $X=1$.
At every finite solver budget, its published schema still admits
$\operatorname{box}(2)$, and its returned caller remains $X=1$.
Replacing the saved birth block by the current block loses that answer.
A separate shared-output producer follows the current caller, admitting
$\operatorname{box}(1)$ and refusing $\operatorname{box}(2)$; its returned
information has a model.

These are output comparisons for actual firings with the specified rule,
selection and output syntax, with concrete solver-backed instances.
The complete selected ordinary/shared body is joined next. General
program-context observations remain separate obligations; invocation
completion is established below.

\begin{theorem}[Conditional constructed-template publication]
A single authored template clause may publish ordinary values, shared
references, and schemas with private holes, explicit actor captures and
latent equations. Both a schema's body and its equations are transported
by the same binding map. Its body decodes from the actual prepared
production, and its conditional instances are exactly the source row's
instances under the returned actor environment.

Every given actor information model extends through the publication, and
every result model projects back. Schema conditions remain local: an
incompatible dormant row does not make the actor information inconsistent.
\end{theorem}
The construction is in \code{ContextualConditionalTemplates},
\code{ContextualConditionalTemplateSource} and
\code{ContextualConditionalTemplateExecution}. It reuses the existing
template production grammar and conditional execution, installing the
authored conditional metadata without a second resource transition.

Resolved clauses, direct templates and compiled native bodies now share
one pinned execution record:
\code{ContextualConditionalExecution.Published}.
It fixes the authored rule, physical selection and actual output code.
The reuse contract explicitly specifies either all flags or the flags at
produced ports. The latter forgets off-output metadata, as the preceding
resolved-clause contract already did. An actual four-output control changes
an off-output flag: every returned row is unchanged, while the exact
flag-function contract rejects the change. Source-specific row decoders
still establish conditional meaning; the common record does not assume it.
Result-frame transport and extension of every given information model are
proved once on the underlying solver-backed firing.

\begin{theorem}[Publication meaning does not factor through the raw firing]
At every finite solver budget, two authored conditional publications can
have the exact same raw firing, including its computed solver return, but
different conditional answer sets. Consequently no observation of that
raw firing alone recovers answer existence for every authored publication.
\end{theorem}
The control swaps only the saved conditions of two schema outputs while
leaving their value code unchanged. The caller establishes $X=0$. The first
row is usable in one publication and unusable in the other. An actual
following clause succeeds for the first and is refused for the second.
These results are in
\code{ContextualConditionalTemplateControls.PublicationBoundary}
and \code{ContextualConditionalTemplateReuse}.

The raw native-step result and the interpreted conditional publication
therefore remain distinct contracts. The constructed publication theorems
retain their stated meaning, but composing the raw native result as though
it already contained the schema conditions is invalid. A canonical program
execution must retain the publication descriptor as execution data, or
encode it in the actual output representation. This boundary is not closed
by the pinned raw-execution record alone.

The direct compiler now retains a structured schema recipe. It contains
an arity-preserving name-role map $\rho$ into private names and captures,
and a finite local equation list $E$, both constructed before execution.
It contains no replacement body. Given the actual emitted term $v$ and
the returned capture assignment $\sigma$, publication computes
\[
  \bigl(\mathsf{schema}(\sigma^{\uparrow}(\rho(v))),
         \sigma^{\uparrow}(\rho(E))\bigr),
\]
where $\sigma^{\uparrow}$ leaves private names fixed. The equations stay
in that row, rather than becoming assertions in the global store.
\code{ContextualConditionalPublication.Recipe} supplies this construction.
\code{ContextualConditionalTemplates.recipe\_publish} proves equality with
the independently constructed conditional template, for both body and
conditions. The actual direct-template target now uses this recipe on its
emitted output. Its row function factors through the raw execution plus
the pre-execution recipes; recipes outside produced ports are absent.
The shared step must preserve this data for each publication family; an
interface exposing only the raw result cannot support this composition.

Mixed native bodies now use the same recipe construction. Their recipes
are computed from surviving authored births before execution, preserving
the schema/shared constructor tag as well as name roles and conditions.
Private names enter the existing tagged result bank; neither the physical
commit nor the global information problem changes. A schema body is decoded
from the actual emitted term, while a live body is that emitted term under
the returned capture assignment.

\begin{theorem}[Native publication from emitted values and recipes]
For every actual compiled mixed-body firing, each selected output row is
exactly the row obtained from the saved source birth, including its
constructor, value and local equations. Under the source shape and resource
conditions, every net output has a recipe, and retained ports have none.
\end{theorem}
The construction and its equality theorem reuse
\code{ContextualConditionalTemplateNativeProducer} and
\code{ContextualConditionalTemplateNativePublicationValues}.
The matching theorem is derived from this row equality. Thus the existing
source-world and invocation comparisons continue to apply to the constructed
result, without a separate resource-query or invocation proof chain.

The packed-bank construction is shared in
\code{ContextualConditionalRowPacking}: the existing execution-private bank
occupies the \code{none} block, and output-private families occupy their
\code{some} blocks.

\code{publishWithFallback} renders recipes from actual
outputs and retains the actual execution rows at ports without recipes.
Its retained-row theorem preserves the complete conditional matching
relation under every future assignment. Mixed native publication is an
instance of this construction.

The concrete mixed clause establishes $X=0$ and publishes two schemas
$\mathsf{pair}(X,\mathsf{pair}(Y,Y))$, conditioned respectively by $X=0$
and $X=1$, beside ordinary and shared outputs. A following actual clause
reads the first schema twice, with private values $0$ and $1$, in one
joint firing. The second schema remains present and unmatched. Selecting
it instead is refused at every solver budget. Repeated $Y$ occurrences
within one use cannot disagree, the captured $X$ cannot be chosen again,
and erasing the condition admits an otherwise invalid source answer.
The controls are in \code{ContextualConditionalTemplateControls} and
\code{ContextualConditionalTemplateReuse}.

The independent update in \code{ContextualConditionalTemplateState}
permits incoming row locals to differ from the clause's query variables.
Its ordinary, shared, captured and retained rows are constructed from
source syntax. \code{ContextualConditionalTemplateClauses} uses the
existing resolved-clause compiler. Its decoder connects actual prepared
outputs to the final published schema body.

\begin{theorem}[Conditional template clause comparison]
At every finite solver budget, source admission holds if and only if an
actual compiled clause exists with the source-equivalent result world.
The comparison preserves physical occurrences, returned actor bindings,
every information model in both directions and every conditional
row-matching observation in fresh declaration extensions.
\end{theorem}
The joined comparison is in
\code{ContextualConditionalTemplateClauseResult}.

\code{ContextualConditionalTemplateSchedules} additionally permits finite
bodies mixing these template clauses with existing ordinary/shared
segments. Birth descriptions retain the authored conditional row and a
separate transport for its captures. Admission compilation preserves and
reflects source admission, using final information, original guard cuts,
port shapes and the complete resource precondition, including intermediate
freshness. The net resource effect gives the exact source occurrence set.
These are admission and resource results; native one-firing realization
for the extended grammar remains open.

The concrete mixed source body publishes the two schemas above, then reads
the compatible occurrence twice with private values $0$ and $1$. Both the
good body and the version selecting the incompatible row have available
physical resources, but only the good body is admitted. Its dormant row
remains present. The actual clause comparison also constructs two
source-linked native firings and excludes the incompatible native second
read, at arbitrary declared budgets. These controls are in
\code{ContextualConditionalTemplateSourceControls}.

\begin{theorem}[Surviving conditional births and native net inputs]
Under the authored port shapes, each final mixed-body row has exactly the
conditional instances of its latest saved production, or of the retained
entry row under the composed capture transport. This comparison holds
before guard or resource admission is established.

Every model of the final source information supplies matching witnesses
for all tracked initial inputs. Any finite list of those inputs can be
opened together in a fresh extension, preserving the given source model's
names. Under source resource executability, every actual allocated net
input is one of these entry occurrences. Every net output selects a saved
conditional birth with exactly the final source row's meaning.
\end{theorem}
The origin comparison is in
\code{ContextualConditionalTemplateOriginSemantics}.

Joint witnesses are in \code{ContextualConditionalTemplateInputModels}.

The resource-net links are in
\code{ContextualConditionalTemplateAllocation}.
The lookup is computed from authored selections and birth descriptions,
without using a final source-row oracle.

The entry control reads a row conditioned by $X=0$ beside an unused row
conditioned by $X=1$. Every source model supplies the used input witness;
the unused row remains incompatible. The allocated net reads only the
used entry. A separate control consumes logical occurrence 10, then
recreates it in a later clause. The new allocation generation is in the
net output set; its old generation and entry name are not. Its meaning
comes from the later production. The entry and logical reuse controls are
in \code{ContextualConditionalTemplateOriginControls}.

The allocation controls are in
\code{ContextualConditionalTemplateAllocationControls}.

\begin{theorem}[Mixed native producer and conditional result]
Assume the mixed body's authored port shapes, complete source resource
executability and consistent final information. At every finite solver
budget there is an actual native binding/resource producer firing with
the constructed rule, selection, output expressions and publication
flags. Its conditional publication has the same visible occurrences,
returned caller bindings and conditional row instances as the source
result, under the declared bound-name projection. Every source model
extends to a returned model, and every returned model projects to a source
model; row instances agree in every fresh future.
\end{theorem}
The mixed body is a compiled body in
\code{ContextualConditionalTemplateNativeProblem}; its binding construction
is the common one.

The actual output decoding is in
\code{ContextualConditionalTemplateNativePublicationValues}.

The world comparison is in
\code{ContextualConditionalTemplateNativeResult}.

Its physical, model and world laws are the common ones of
\code{ContextualNativeResults}.

Output-private declarations carry no global assertion of their row
conditions. Value and conditional metadata use the same role transport.
The controls fuse two compatible private reads, preserve an unused
incompatible entry, and exclude a firing that selects the incompatible
row.

The controls are in
\code{ContextualConditionalTemplateNativeControls}.

This is a result comparison; admission needs the original-cut data
queries. A control has consistent information and available resources
while a private creation refutes a later absence guard. The source is
refused, but the binding/resource producer executes. Thus producer
existence alone cannot replace source admission; with the queries
attached, the same body has no admitted firing. This control is in
\code{ContextualConditionalTemplateNativeGuardBoundaryControls}.

Every original cut of a mixed body is represented over entry occurrences
and authored prefix creations, with the same problem, checks and guard
meaning. The joined admission theorem below is therefore stated once for
compiled bodies.

The mixed cut representation is in
\code{ContextualConditionalTemplateEntryCuts}.

\begin{theorem}[Selected private body to one native firing]
For a finite authored body, with ordinary and shared productions and with
selectively captured conditional templates, source admission holds if and only if there exists one
actual query-admitted native firing with the constructed rule, physical
selection, output function, publication flags and query list, at any
declared finite solver budget. No source shape, resource-validity or
successful-execution premise is supplied to the native side.

The invocation control is consumed. After decoding bound creation names,
its world is conditionally equivalent to the source result:
physical occurrences agree, every source information model extends to a
native model, every native model projects to a source model, returned
bindings agree, and conditional row-matching observations agree in every
fresh declaration extension.
\end{theorem}
The admission theorem is in \code{ContextualNativeProgramExecution}, for
every compiled body and producer.

The joined theorem is \code{source\_iff\_native\_body}.

For ordinary/shared bodies it and the world comparison are in
\code{ContextualNativeProgramResult}; for mixed bodies they are in
\code{ContextualConditionalTemplateNativeResult}, where captured rows keep
their conditions as row metadata.

The producer reads authored birth descriptors, not the final source row
store. Its ordinary values retain their birth problem; shared values use
the returned caller. Resource queries check required and forbidden entry
identities against actual inventory data. The original finite compiled
resource precondition includes intermediate freshness tests that the net
effect alone can erase. Static port shapes and its store-independent
possibility flag are checked separately on the control row. Guard queries
then retain each clause's original information cut, including its own
port and tell equations.

The concrete compiler instances include a stored-lambda application that
computes a lookup key, takes a permit and reuses its answer twice; an
incompatible reuse is refused. An earlier reusable schema keeps its
answers, whereas a later tell cannot repair an earlier failed guard.
A count guard uses its own clause's equation. A make-then-take program
with an occupied creation identity has an enabled net effect but fails
the complete resource condition and its actual native absence query;
the same internal program is admitted from an empty entry.

The comparison fixes the selected authored body and uses the stated
conditional row-and-model observation. Its entry control and helper
encodings remain observable to unrestricted native programs.
Arbitrary target selections and program observers, full active evaluation, collection
composition, fairness and cost retain their separate obligations.
Semantic support tables do not supply a total higher-order matching or
guard-decision procedure.

\begin{theorem}[Conditional preparation and private publication transport]
Mapping captured declarations commutes with the complete firing-information
store: incoming equations, matched values, authored equations and each
selected row's latent conditions. Take/read scopes and the physical world
are unchanged. The corresponding original-cut guard ask is preserved.
This does not assert that every capture map preserves solvability: the
control substituting zero for a captured value constrained to one is
inconsistent after transport.

An arity-preserving reindexing of private row names commutes with publication
and with capture transport, preserving the schema/shared constructor, value
and pending equations. Reindexing cannot invent a matching answer; an
arity-preserving left inverse also gives the converse. Thus private
renaming preserves matching without identifying captured references.
\end{theorem}
The preparation and guard laws are in
\code{ContextualConditionalExecution}; the private-row and recipe laws are
in \code{ContextualConditionalPublication}. The actual pending-answer
control still accepts one and rejects two after embedding its private bank
in a fresh tagged bank. The mixed native consumer checks that reindexing
its compiled recipes agrees exactly with reindexing its published rows.
These are transport components. They do not establish the full native
commuting diamond or equivalence for all following clients.

\begin{theorem}[Relocation of actual conditional firings]
Let a permutation align every freshly allocated rule, take and read scope
between two occupied-scope inventories, fixing the captured summand.
Relocating the output expressions preserves and reflects conditional
admission, including the full heterogeneous information and original-cut
count and absence asks. Every admitted relocated operation has an actual
solver return at every finite budget, with the same physical occurrence
effects. Ordinary reusable outputs have exactly the same matching instances
in both executions, even when the returned helper families differ and
residual equations remain.
\end{theorem}
The preparation laws are in \code{ContextualEventViewRelocation}.

Conditional asks and source instances are handled by
\code{ContextualConditionalNaturality}.

Admission and firing relocation are in
\code{ContextualConditionalExecution}. Adding scope reservations preserves
formation in \code{ContextualConditionalScopes}.

A concrete relocated
consumer takes a pending answer, checks count and absence, and still
publishes seven; an incompatible take remains refused. The execution
component of an actual compiled mixed body also consumes the relocation
theorem. Transport of its native predicate-query list and selectively
captured publication recipes remains separate, as do private-family changes
after intervening effects and the full two-firing diamond.

\begin{theorem}[Selected invocation completion]
The compiled firing takes its invocation control after the data takes
and never produces a replacement control. Every successful invocation
therefore removes that control and strictly decreases its pending
indicator from one to zero, while preserving the source/result comparison.
No subsequent admission requiring the same control is possible after any
intervening data steps that do not produce a control occurrence.
\end{theorem}
The completion facts are in \code{ContextualNativeInvocation}, for every
compiled body; the strengthened source comparisons are
\code{source\_iff\_completed\_invocation} for ordinary/shared and for mixed
bodies.
Appending the control preserves every data opening ordinal and the
control's own former read ordinal. This is ordinary resource consumption,
not a new identity or occurrence-lifetime model. A fresh control production
is an explicit boundary; standing installation, cancellation and retry
contracts remain separate.

The concrete empty source body has an actual native firing at every finite
solver budget. Its public result is empty, its private control is consumed,
and no further invocation admission is possible without a new control.
By contrast, a read gate is enabled for arbitrarily many resource
executions without changing its store. The completion result does not
infer successful execution merely from control absence in an arbitrary
environment.

\section{From interaction to the full language}

The following comparisons remain the acceptance criteria for a common
foundation. Existing restricted results are inputs to them.
\begin{center}
\begin{tabularx}{\linewidth}{@{}lX@{}}
\toprule
Construction & Required account \\
\midrule
Querying & Complete selection, answer multiplicity, bindings and any required order. \\
Communication & Consumption or shared reading, continuation activation and contextual name equality. \\
Expression execution & Application, lexical binding, shared versus independent choice, effects and suspension. \\
Reflection & Inert quoted data, explicit activation and declared restrictions on syntax observations. \\
Changing definitions & The equations and versions used by each computation and retained conclusion. \\
Checked inference & Ordinary rule construction corresponding to independent inference judgments, with hypotheses and authority retained. \\
\bottomrule
\end{tabularx}
\end{center}

The target is a concise semantic account plus qualified constructions and
obstructions. Runtime convenience operations, source filenames, and
module counts do not establish these comparisons. Internal administrative
steps need a termination or explicit suspension argument; costs and
fairness have their own assumptions. The candidate is retained for this
investigation, with the full foundation judgment still open.

\begin{thebibliography}{9}
\bibitem{logic-distributive}
M. Stay and L. G. Meredith.
\emph{Logic as a distributive law}. 2016.
\url{https://arxiv.org/abs/1610.02247}.
\bibitem{metametta}
L. G. Meredith, B. Goertzel, J. Warrell and A. Vandervorst.
\emph{Meta-MeTTa: an operational semantics for MeTTa}. 2023.
\url{https://arxiv.org/abs/2305.17218}.
\bibitem{dcalc}
B. Goertzel. \emph{The Distinction Calculus: d-Forms for Hives of
Reflexively Self-Modifying Systems}. Corrected manuscript, September 2026.
\bibitem{hyperseed}
\emph{Hyperseed in the d-Calculus}, mathematical companion to Ben Goertzel's
\emph{Hyperseed}, Chapter 7: signed event calculus and explicit absence.
Version 1, September 2026.
\bibitem{lybech}
S. Lybech. \emph{Encodability and Separation for a Reflective Higher-Order
Calculus}. EPTCS 368, 2022, pp.\ 95--112, Section 2.
\url{https://arxiv.org/abs/2209.02356}.
\end{thebibliography}
\end{document}
